% Options for packages loaded elsewhere
\PassOptionsToPackage{unicode}{hyperref}
\PassOptionsToPackage{hyphens}{url}
%
\documentclass[
]{article}
\usepackage{amsmath,amssymb}
\usepackage{iftex}
\ifPDFTeX
  \usepackage[T1]{fontenc}
  \usepackage[utf8]{inputenc}
  \usepackage{textcomp} % provide euro and other symbols
\else % if luatex or xetex
  \usepackage{unicode-math} % this also loads fontspec
  \defaultfontfeatures{Scale=MatchLowercase}
  \defaultfontfeatures[\rmfamily]{Ligatures=TeX,Scale=1}
\fi
\usepackage{lmodern}
\ifPDFTeX\else
  % xetex/luatex font selection
\fi
% Use upquote if available, for straight quotes in verbatim environments
\IfFileExists{upquote.sty}{\usepackage{upquote}}{}
\IfFileExists{microtype.sty}{% use microtype if available
  \usepackage[]{microtype}
  \UseMicrotypeSet[protrusion]{basicmath} % disable protrusion for tt fonts
}{}
\makeatletter
\@ifundefined{KOMAClassName}{% if non-KOMA class
  \IfFileExists{parskip.sty}{%
    \usepackage{parskip}
  }{% else
    \setlength{\parindent}{0pt}
    \setlength{\parskip}{6pt plus 2pt minus 1pt}}
}{% if KOMA class
  \KOMAoptions{parskip=half}}
\makeatother
\usepackage{xcolor}
\usepackage[margin=1in]{geometry}
\usepackage{graphicx}
\makeatletter
\def\maxwidth{\ifdim\Gin@nat@width>\linewidth\linewidth\else\Gin@nat@width\fi}
\def\maxheight{\ifdim\Gin@nat@height>\textheight\textheight\else\Gin@nat@height\fi}
\makeatother
% Scale images if necessary, so that they will not overflow the page
% margins by default, and it is still possible to overwrite the defaults
% using explicit options in \includegraphics[width, height, ...]{}
\setkeys{Gin}{width=\maxwidth,height=\maxheight,keepaspectratio}
% Set default figure placement to htbp
\makeatletter
\def\fps@figure{htbp}
\makeatother
\setlength{\emergencystretch}{3em} % prevent overfull lines
\providecommand{\tightlist}{%
  \setlength{\itemsep}{0pt}\setlength{\parskip}{0pt}}
\setcounter{secnumdepth}{-\maxdimen} % remove section numbering
% definitions for citeproc citations
\NewDocumentCommand\citeproctext{}{}
\NewDocumentCommand\citeproc{mm}{%
  \begingroup\def\citeproctext{#2}\cite{#1}\endgroup}
\makeatletter
 % allow citations to break across lines
 \let\@cite@ofmt\@firstofone
 % avoid brackets around text for \cite:
 \def\@biblabel#1{}
 \def\@cite#1#2{{#1\if@tempswa , #2\fi}}
\makeatother
\newlength{\cslhangindent}
\setlength{\cslhangindent}{1.5em}
\newlength{\csllabelwidth}
\setlength{\csllabelwidth}{3em}
\newenvironment{CSLReferences}[2] % #1 hanging-indent, #2 entry-spacing
 {\begin{list}{}{%
  \setlength{\itemindent}{0pt}
  \setlength{\leftmargin}{0pt}
  \setlength{\parsep}{0pt}
  % turn on hanging indent if param 1 is 1
  \ifodd #1
   \setlength{\leftmargin}{\cslhangindent}
   \setlength{\itemindent}{-1\cslhangindent}
  \fi
  % set entry spacing
  \setlength{\itemsep}{#2\baselineskip}}}
 {\end{list}}
\usepackage{calc}
\newcommand{\CSLBlock}[1]{\hfill\break\parbox[t]{\linewidth}{\strut\ignorespaces#1\strut}}
\newcommand{\CSLLeftMargin}[1]{\parbox[t]{\csllabelwidth}{\strut#1\strut}}
\newcommand{\CSLRightInline}[1]{\parbox[t]{\linewidth - \csllabelwidth}{\strut#1\strut}}
\newcommand{\CSLIndent}[1]{\hspace{\cslhangindent}#1}
\usepackage{booktabs}
\usepackage{longtable}
\usepackage{array}
\usepackage{multirow}
\usepackage{wrapfig}
\usepackage{float}
\usepackage{colortbl}
\usepackage{pdflscape}
\usepackage{tabu}
\usepackage{threeparttable}
\usepackage{threeparttablex}
\usepackage[normalem]{ulem}
\usepackage{makecell}
\usepackage{xcolor}
\ifLuaTeX
  \usepackage{selnolig}  % disable illegal ligatures
\fi
\usepackage{bookmark}
\IfFileExists{xurl.sty}{\usepackage{xurl}}{} % add URL line breaks if available
\urlstyle{same}
\hypersetup{
  hidelinks,
  pdfcreator={LaTeX via pandoc}}

\author{}
\date{\vspace{-2.5em}}

\begin{document}

\section{Riverine Heatwaves}\label{rh}

\emph{Author: Tom Baumann}

\emph{Supervisor: Henri Funk}

\emph{Degree: Bachelor}

\subsection{Abstract}\label{abstract}

Riverine Heatwaves, which are prolonged periods of exceptionally high
water temperatures in river systems, are a hydrological phenomenon, that
is expected to gain considerable importance throughout the next years,
also due to their link to climate change. The following paper will give
an introduction to the precise definition of riverine heatwaves, as well
as explain its drivers and impacts. Furthermore, it will focus on the
methodology used to detect riverine heatwaves from temperature time
series, and how their development, both spatially and temporally, can be
quantified. Applying these methods on a dataset from Main river system
in Bavaria delivered a significant increase in frequency over time,
while increasing temporal trends in duration and intensity were
insignificant. Also linear trends between river kilometer and the
heatwave metrics could not be verified. However, the dataset only
included a relatively short time period, therefore a larger data size
might lead to more precise results.

\subsection{Introduction}\label{introduction}

\subsubsection{Overview}\label{overview}

In contrast to air, lake and ocean heatwaves, research on river
heatwaves is still just emerging. However it is a topic that is only
projected to become more relevant over time, as increasing water
temperatures in rivers are closely linked to climate change, and also
riverine heatwaves are expected to become more frequent, long and
intense. In fact, they are even projected to increase more rapidly than
air heatwaves (Sadayappan \& Li, 2025). To research riverine heatwaves,
it is important to understand precisely how they are defined, how they
can be quantified, and what sort of data is required for their analysis.
Also the drivers and impacts of riverine heatwaves need to be examined,
to understand their appearance, and how they can be mitigated. Lastly,
it is important to be able to forecast and predict riverine heatwaves,
which requires the construction of reliable models. This paper is going
to give an overview on the topic of riverine heatwaves, by introducing
their definition, drivers and impacts. This will be helpful, to gain a
better understanding of riverine heatwaves. Furthermore it will give a
baseline on how riverine heatwaves can be detected from temperature time
series, and what metrics can be used to quantify them. Later these
methods will be used to perform an exploratory analysis on a data set of
water temperature data from Main river system, focusing on the evolution
of riverine heatwaves over time and the differences between stations.

\subsubsection{Definition}\label{definition}

Riverine Heatwaves are defined similarly to atmospheric heatwaves, as
periods where daily mean water temperature exceeds a local seasonally
varying percentile based threshold for a prolonged period of time. For
riverine heatwaves, usually a 90th percentile threshold and a 5 day
period are used (Hamel et al., 2025). Although their definitions are
similar, riverine heatwaves and atmospheric heatwaves behave
differently, as riverine heatwaves tend to be more rare, but more
persistent, than atmospheric heatwaves (Sadayappan \& Li, 2025).

\subsubsection{Drivers}\label{drivers}

As riverine heatwaves are a phenomenon that has not gained much
attention so far, its drivers are not fully understood yet. We know that
river temperature in general is controlled by a variety of factors.
There are the atmospheric conditions like air temperature, wind,
rainfall and solar radiation, the hydrologic conditions which means the
discharge, how the river interacts with different water bodies within
its proximity and processes at the water stream bed interface, for
example the interaction of the river with ground water. On top of that,
the water temperature can be influenced by its environment, for example
by the amount of vegetation and shading, by impervious surfaces nearby
or by the presence of thermal effluents from near industries or human
waste water (Hamel et al., 2025).

However, so far little is known on how specifically riverine heatwaves
are caused. From literature we know, that riverine heatwaves are
primarely driven by climate drivers. These include air temperature,
while especially trends of minimum and maximum temperature play a role,
and trends in low discharge. Also the contribution of groundwater,
winter precipitation percentage and trends of snow water play a role, as
these help to buffer the impact of extremely hot air temperatures on
river temperature (Sadayappan \& Li, 2025). Behind climate drivers,
there are also direct anthropogenic drivers, such as river dams, which
typically tend to elongate riverine heatwaves (Sadayappan \& Li, 2025).
Nevertheless, the processes that cause riverine heatwaves are still not
fully understood.

\subsubsection{Impacts}\label{impacts}

The appearance of riverine heatwaves can have both ecological and
socioeconomic impacts. These can be parted into different layers.
Firstly there are the immediate effects of prolonged periods of high
water temperature. The higher temperature could lead to an increased
rate of open water evaporation. Moreover, they might lead to an
increased spread of diseases, harmful algee bloom and a decrease in
water quality in general. This in turn can reduce the availability of
drinking water or the possibility of recreational activities and be
harmful to flora and fauna living in, or close by the river systems.
Also the prolonged high temperatures can cause thermal stress on fish
and other aquatic species (Hamel et al., 2025). Together this is
threatening fishing and aquaculture which is a vital source of food for
a large part of the population (Chen et al., 2026). Moreover, high water
temperatures can also negatively impact energy production, as nuclear
power production often relies on cooling water from close by river
system, which already has been unavailable due to riverine heatwaves in
some cases (Hamel et al., 2025).

\subsection{Data}\label{data}

\subsubsection{Data Set}\label{data-set}

For the analysis a dataset of longitudinal temperature data from 14
gauges at Main river system in Bavaria was used. The data set includes
eight measurements of water temperature per day over a time period of 5
to 20 years between 2000 and 2020. According to the literature, the
analysis of riverine heatwaves requires time series of at least 10,
ideally 30 years with at least daily data and not more than 50 percent,
ideally not more than 25 percent of missing data (Hamel et al., 2025) .
Also for this analysis it is necessary to have a common time period for
all stations, to obtain comparable results between stations. As many of
the stations do not cover a common time period, and as some of the
stations have a high amount of missing data for some years, the analysis
was limited to six of the stations to guarantee these standards.
Excluding the other stations gave us a 15 year common period time series
of daily water temperature, with almost no missing values (not more than
ten percent per yearly station specific time series). The used stations
are marked in figure @ref(fig:rh\_stations-map), which also shows their
position along the main river system.

\subsubsection{Locations}\label{locations}

\begin{figure}
\centering
\includegraphics{05-riverine_heat_files/work/05-riverine_heat/figures/rh_05-stations_map.png}
\caption{Map of Main river systen. The station included in the analysis
are marked red.}
\end{figure}

\subsection{Methods}\label{methods}

\subsubsection{Threshold}\label{threshold}

The threshold which classifies riverine heatwaves has to be precisely
defined. This paper uses the definition of marine heatwaves by Hobday.
Although it was originally defined for marine heatwaves, it has since
been used on many different studies about riverine heatwaves as well.

The definition of marine heatwaves by Hobday uses a 90th percentile
seasonally varying threshold, which is also varying with location and
therefore calculated separately for each station. To calculate the
seasonally varying threshold, for each day the historical temperature
measurements, as well as the historical measurements from a 11 day
surrounding window are taken. From these data points, the 90 percent
quantile is taken. The exact formula for this calculation is depicted
below, where y is the year variable, d is the day variable and j is the
day for the considered temperature @ref(eq:rh\_threshold). \(Y_s\) and
\(y_e\) are the years that mark the beginning and end of the baseline
period, \(P_{90}(X)\) is the 90 percent quantile of \(X\) . This
procedure will be repeated for each day of year, until the threshold for
the entire year is completed (Hobday et al., 2016).

\begin{equation}
T_{90}(j) = P_{90}(X),
\quad \text{where}
\quad
X = \left\{T(y,d)\;\middle|\; y_s \le y \le y_e,\; j-5 \le d \le j+5 \right\}
(\#eq:rh_threshold)
\end{equation}

Usually, for this a historical baseline period is used. But since here
the data set only covers a relatively short time period, the same time
period was used for the calculation of the threshold and the analysis.
Even if the same threshold definition is used, comparing results from
different studies on riverine heatwaves has to be treated carefully. It
can be problematic to compare exact numbers for metrics between studies
with different baseline periods, as climate change suggests a
temperature biased over time. However, as in this paper the focus is
only to demonstrate how heatwaves evolve relatively to another over time
by looking for general trends rather than comparing exact metrics with
other studies. And since in this paper only measurements from the same
stations with the same baseline period were compared, this approach also
suffices.

Resulting from this procedure, for each station a threshold is obtained,
that can be used for comparison with the time series of each year. In
figure @ref(fig:rh\_threshold-kemmern) an example of this is shown for
the Kemmern station and the year 2018. Heatwaves can then be detected,
by firstly marking all potential candidates, which are all days
exceeding the thresholds. After, the time series are examined for events
of at least five consecutive days exceeding the threshold. These are
marked as heatwaves and given an ID, while all shorter events are
discarded. When there is a gap of only one or two days between two or
more heatwaves that have a duration of at least five days each, the
heatwaves, are merged into one heatwave.

\subsubsection{Threshold example: Main data (Kemmern,
2018)}\label{threshold-example-main-data-kemmern-2018}

\begin{figure}
\centering
\includegraphics{05-riverine_heat_files/work/05-riverine_heat/figures/rh_05-threshold_kemmern.png}
\caption{Temperature time series for Main Kemmern in 2018 (black line).
The dotted orage line represents the heatwave threshold for Kemmern
station. The red areas mark heatwaves in that year for Kemmern station.
The light grey lines represent all temperature time series of Kemmern
station for the baseline period.}
\end{figure}

\subsubsection{Metrics used}\label{metrics-used}

This paper focuses on three primary metrics, that are used to quantify
the evolution of riverine heatwaves. A frequency metric, a duration
metric and an intensity metric were used, to describe the evolution
under different aspects.

To quantify the evolution of frequency the number of heatwave events per
year is used. For this, the number of different heatwaves IDs are being
counted per year and station.

Furthermore, the duration of a heatwave is considered. This is
calculated using the first and last day of each heatwave
@ref(eq:rh\_duration) (Hobday et al., 2016). The duration plays an
important role, as especially longer heatwaves can have more severe
impacts.

\begin{equation}
D = t_e - t_s
(\#eq:rh_duration)
\end{equation}

Lastly, intensity is used, to assess the amount of the heat spike. The
original definition by Hobday suggests to to define it as the
temperature by which the climatological mean is exceeded. However, most
literature on riverine heatwaves define intensity as the exceendance of
the 90th percentile threshold (Sadayappan \& Li, 2025) (Chen et al.,
2026). Therefore this was also adopted for this paper. Also in this
paper intensity is defined, as the mean of the daily temperature
exceedance per heatwave@ref(eq:rh\_intensity).

\begin{equation}
i_{\mathrm{mean}}
=
\overline{T(t)-T_{90}(j)}
(\#eq:rh_intensity)
\end{equation}

These metrics can also be used to summarize two or more of these
characteristics. For example, the heatwave events can be multiplied with
the mean annual duration, to obtain the number of heatwave days per
year. Or the duration of a heatwave can be multiplied with its mean
intensity, to obtain its mean severity.

\subsubsection{Research Question}\label{research-question}

This paper examines the evolution of riverine heatwaves during climate
change. Therefore, its focus lies primarily on the development of the
three previously defined metrics over time by posing the following
questions:

\begin{enumerate}
\def\labelenumi{\arabic{enumi}.}
\item
  How has the frequency of riverine heatwaves changed over time?
\item
  How has the duration of riverine heatwaves changed over time?
\item
  How has the intensity of riverine heatwaves changed over time?
\end{enumerate}

Given the availability of stations with different locations along the
same river system, it was also looked into how the metrics differ
between the stations. To answer these research questions, an exploratory
analysis was performed on the heatwave data, examining it for
differences between stations and years. Also linear and quasi poisson
modells were fitted on the data, to look for significant temporal and
spatial trends.

\subsection{Results}\label{results}

\subsubsection{Frequency}\label{frequency}

Figure @ref(fig:rh\_heatwave-events-overview) shows the distribution of
the number of heatwave events per year for each station, with the
stations being ordered by their position along the river from downstream
to upstream. The plot shows, that for each station the mean number of
heatwave events lies between 2.5 and 3. Both its mean and its variance
seem to slightly increase in a downstream direction. To investigate this
more thoroughly, a linear model between the mean number of heatwave
events per station, and the river kilometer was used. The river
kilometer is a static variable that describes, the stations distance to
the river mouth in kilometers. The model indicated a decreasing, however
insignificant trend of -0.001935 events / year with a p value of 0.197.
As Pettstadt station lies on an inflow river and not directly on Main,
it was excluded from this model.

To examine the temporal trend of frequency, the heatwave events were
grouped by their station, and a quasi poisson model was used, to look
for correlation between the frequency and year. Fitting a model for each
individual station led only to 1 significant and 5 insignificant
results. Furthermore a common temporal trend quasi poisson model of
frequency and year was used, to examine the general temporal frequency
trend. An analysis-of-deviance F-test was used to assess whether
including a year x station interaction significantly improved the quasi
poisson model compared with a model containing only year and station
main effects. The year × station interaction did not significantly
improve model fit (analysis of deviance F-test: F(5, 78) = 0.737, p =
0.598), indicating no statistical evidence that temporal trends differed
among stations.Therefore, the model with station as a fixed effect and a
common temporal trend was used for the analysis. This resulted into a
significant positive trend over time (slope: 0.05350, p value : 0.002),
which indicates an increase of heatwave events by 5.5 \% per year. The
common quasi poisson is depicted in figure
@ref(fig:rh\_heatwave-events-plot) where it is plotted into the
frequency distribution of each station individually. The offset
represents the intercept of the station effect, as the frequency of
heatwaves differs from station to station.

\begin{figure}
\centering
\includegraphics{05-riverine_heat_files/work/05-riverine_heat/figures/rh_05-heatwave_events_overview.png}
\caption{Frequency distribution over all stations. The orange square
signifies station mean value.}
\end{figure}

\begin{figure}
\centering
\includegraphics{05-riverine_heat_files/work/05-riverine_heat/figures/rh_05-heatwave_events_plot.png}
\caption{Temporal development of heatwave events by station. The blue
line shows the common temporal trend quasi poisson model.}
\end{figure}

\subsubsection{Duration}\label{duration}

The durations of heatwaves of Main river in the considered period lay
between 5 and 28, with averages between 8 and 10 days. Again the
stations were placed according to their order along the river, and while
the mean stays consistent, the variance seems to increase downstream
@ref(fig:rh\_mean-duration-overview). Also extremely long heatwaves seem
more common at the downstream stations. Again there was a model fitted
between the mean duration of each station and the river kilometer, this
time a linear model was used, however it showed no significant trend
(slope: -0.002058 days / km, p value: 0.459). Another idea would be to
examine the mean of the yearly maximum duration of each station and look
for a downstream trend. However, for this data set, this approach also
did not deliver a significant result (slope: -0.002475 days / km, p
value: 0.616 ).

For the temporal trend a similar approach was used. The yearly mean
durations for each station were plotted @ref(fig:rh\_duration-plot), and
for each a linear model was fitted. Four of the six stations showed a
positive trend, with only one of them being significant, while the other
stations showed a insignificant negative trend. To fit a linear model
for the common temporal duration trend, again the model with year x
station interaction was tested against the model including only the main
effects, analogously to the frequency trend, this time using a analysis
of variance F-Test. Again there was no indication for significantly
differing trends between station, therefore the model with only the main
effects was chosen. This resulted in a increasing, however
insignificant, temporal trend of duration (slope: 0.08490 days / year, p
value: 0.217).

\begin{figure}
\centering
\includegraphics{05-riverine_heat_files/work/05-riverine_heat/figures/rh_05-mean_duration_overview.png}
\caption{Distribution of mean duration for all stations. Each point
signifies a single heatwave. The orange square signifies the average
duration per station.}
\end{figure}

\begin{figure}
\centering
\includegraphics{05-riverine_heat_files/work/05-riverine_heat/figures/rh_05-duration_plot.png}
\caption{Temporal development of mean heatwave duration by station. Each
point signifies the annual mean duration for the considered station. The
blue line shows the common temporal trend linear model.}
\end{figure}

\subsubsection{Intensity}\label{intensity}

The intensity lies between 0 and 2 degrees celsius over the threshold,
with the stations means laying between 0.5 and 0.8 degrees Celsius.
Again the plot @ref(fig:rh\_mean-intensity-overview) suggests a trend
between the location of the stations, it seems like the mean intensity
and its variance decrease stream downwards. To capture this trend, a
linear model was used, which gave a slope of 0.0011531 °C per km,
however again this result was not significant for the 95 percent
confidence level with a p value of 0.059.

Looking at the yearly data @ref(fig:rh\_intensity-plot), the plots again
reflect the higher variance for upstream stations Schwürbitz,
Schweinfurt and Kemmern, however a clear temporal trend is not visible.
Also neither the station specific linear models, nor the common temporal
trend model showed a significant trend. The result of the common model
was a slope 0.008265 °C per year with a p value of 0.187, while the
individual models showed no significant trend for any of the stations.
The common temporal trend model was fitted and tested analogously to the
model for duration.

\begin{figure}
\centering
\includegraphics{05-riverine_heat_files/work/05-riverine_heat/figures/rh_05-mean_intensity_overview.png}
\caption{Distribution of heatwave intensity by station. Each point
represents the intensity of one heatwave. The orange square signifies
the mean heatwave intensity for each station}
\end{figure}

\begin{figure}
\centering
\includegraphics{05-riverine_heat_files/work/05-riverine_heat/figures/rh_05-intensity_plot.png}
\caption{Temporal development of mean heatwave intensity by station.
Each point signifies the annual mean heatwave intensity for the
considered station. The blue line shows the common temporal trend linear
model.}
\end{figure}

\subsection{Discussion}\label{discussion}

\subsubsection{Overview of trends common
models}\label{overview-of-trends-common-models}

\begin{figure}
\centering
\includegraphics{05-riverine_heat_files/work/05-riverine_heat/figures/rh_05-polygon_temporal.png}
\caption{All three heatwave metrics by year. Values are expressed
relative to the mean across years (1.0 = overall mean); the dashed grey
triangle indicates the overall mean.}
\end{figure}

The observed results of all three metrics are summarized for the
temporal trends in figure @ref(fig:rh\_polygon-temporal) and for the
differences between stations in figure @ref(fig:rh\_polygon-spatial).
From the literature a increasing trend over time for all three metrics
would be expected (Chen et al., 2026) (Sadayappan \& Li, 2025). This
analysis only reflects a significant increasing trend for frequency,
while duration and intensity also show increasing, but insignificant
trends. However, the recommended data size of time series with a length
of ideally at least 30 years were not met, as only a period of 15 years
was considered. Also there was no historical period used to calculate
the heatwave threshold. This might explain the insignificance of the
duration and intensity trend. Moreover, only data from one river was
considered, and a trend of one river does not necessarily match the
global trend of riverine heatwave increase. With a longer time series it
might be possible to obtain significant results for the temporal trends
of duration and intensity, or it might be possible to have a larger
confidence that there is no significant temporal trend for these metrics
in Main river system.

However, finding usable data for the analysis of riverine heatwaves is
difficult, as due to the suspected climate change bias it is not always
reasonable to use a long period of historical data as a baseline. With
the temperatures being in general a lot higher than in the past, a
historical threshold would not be able to capture heatwaves anymore, as
a too large amount of temperatures would exceed it. Leading to permanent
heatwave states, which would make the examination of frequency and
duration meaningless, and a temporal trend would not be able to be
detected. An idea to account for the temporal bias in temperature data
would be to use simulated data. By using simulated data, longer time
series could be used, the climate change bias would be incorporated in
the threshold calculation, and predictions for future development of
riverine heatwaves would be possible.

\begin{figure}
\centering
\includegraphics{05-riverine_heat_files/work/05-riverine_heat/figures/rh_05-polygon_spatial.png}
\caption{All three heatwave metrics by station. Values are expressed
relative to the mean across all stations (1.0 = overall mean); the
dashed grey triangle indicates the overall mean.}
\end{figure}

The plots indicated, that there was a correlation between frequency and
intensity and the river kilometer. The implication was, that the
frequency of heatwaves increase downstream ward, while their intensity
decreases. These trends however could not be verified using linear
models, which did not deliver significant results. However, it could be
interesting to perfom further examinations on this trends, using longer
time series, or a larger number of stations along the same river, maybe
also on a different river with a denser network of gauges. The frequency
trend might indicate, that there is a downstream movement of riverine
heatwaves, that might explain the higher frequency downstream wards.
This could be interesting to make predictions on the appearance of
heatwaves downstreams, after one has appeared further upstream.

The results, in a way, show the difficulties of data scarcities when
working on temporal trend analysis with riverine heatwaves during
climate change. This underlines the importance of constructing a dense
network of temperature data, perhaps with the help of simulated data.
Although a significant temporal trend was only detected for frequency,
there were some implications for increasing trends from the other
metrics, which, as potential trends between the river kilometer and
frequency and intensity, might be an interesting relation to examine in
a larger data set. In conclusion riverine heatwaves remain a relevant
phenomenon, during climate change, that still offers many interesting
research possibilities.

\phantomsection\label{refs}
\begin{CSLReferences}{1}{0}
\bibitem[\citeproctext]{ref-chen2026}
Chen, Y., Su, Z., Woolway, R. I., Wanders, N., Wu, S., Huang, Z., \&
Luo, M. (2026). Persistent river heatwaves are emerging worldwide under
climate change. \emph{Nature Communications}, \emph{17}, 94.
\url{https://doi.org/10.1038/s41467-025-66868-5}

\bibitem[\citeproctext]{ref-vanhamel2025}
Hamel, A. van, Bruno, G., Chartier-Rescan, C., et al. (2025). Riverine
heatwaves are an emergent climate change risk. \emph{Nature Water},
\emph{3}, 1356--1364. \url{https://doi.org/10.1038/s44221-025-00541-5}

\bibitem[\citeproctext]{ref-hobday2016}
Hobday, A. J., Alexander, L. V., Perkins, S. E., Smale, D. A., Straub,
S. C., Oliver, E. C. J., Benthuysen, J. A., Burrows, M. T., Donat, M.
G., Feng, M., Holbrook, N. J., Moore, P. J., Scannell, H. A., Sen Gupta,
A., \& Wernberg, T. (2016). A hierarchical approach to defining marine
heatwaves. \emph{Progress in Oceanography}, \emph{141}, 227--238.
\url{https://doi.org/10.1016/j.pocean.2015.12.014}

\bibitem[\citeproctext]{ref-sadayappan2025}
Sadayappan, K., \& Li, L. (2025). Riverine heat waves on the rise,
outpacing air heat waves. \emph{Proceedings of the National Academy of
Sciences}, \emph{122}(39), e2503160122.
\url{https://doi.org/10.1073/pnas.2503160122}

\end{CSLReferences}

\end{document}
